% Part: second-order-logic
% Chapter: metatheory
% Section: second-order-arithmetic

\documentclass[../../../include/open-logic-section]{subfiles}

\begin{document}

\olfileid[pt]{sol}{met}{spa}

\olsection{Aritmética de segunda ordem}

Lembre-se de que a teoria $\Th{PA}$ da aritmética de Peano contém os oito
axiomas de~$\Th{Q}$,
\begin{align*}
& \lforall[x][\eq/[x'][\Obj 0]]\\
& \lforall[x][\lforall[y][(\eq[x'][y'] \lif \eq[x][y])]]\\
& \lforall[x][(\eq[x][\Obj 0] \lor \lexists[y][\eq[x][y']])]\\ 
& \lforall[x][\eq[(x + \Obj 0)][x]]\\
& \lforall[x][\lforall[y][\eq[(x + y')][(x + y)']]]\\
& \lforall[x][\eq[(x \times \Obj 0)][\Obj 0]]\\
& \lforall[x][\lforall[y][\eq[(x \times y')][((x \times y) + x)]]]\\
& \lforall[x][\lforall[y][(x < y \liff \lexists[z][\eq[(z' + x)][y]])]]\\
\intertext{além de todas as sentenças da forma}
& (!A(\Obj 0) \land \lforall[x][(!A(x) \lif !A(x'))]) \lif \lforall[x][!A(x)].
\end{align*}
Esta última é um ``esquema'', isto é, um padrão que gera infinitas
!!{sentence}s da linguagem da aritmética, uma para cada !!{formula}~$!A(x)$.
Chamamos esse esquema de \emph{esquema de axiomas de indução} (de primeira
ordem). Na aritmética de Peano de \emph{segunda ordem}~$\Th{PA^2}$, a
indução pode ser enunciada por uma única sentença. $\Th{PA^2}$ consiste nos
oito primeiros axiomas acima mais o \emph{axioma de indução} (de segunda
ordem):
\[
\lforall[X][((X(\Obj 0) \land \lforall[x][(X(x) \lif X(x'))]) \lif \lforall[x][X(x)])].
\]
Ele diz que, se um subconjunto~$X$ do !!{domain} contém
~$\Assign{\Obj{0}}{M}$ e, com cada~$x \in \Domain{M}$, também contém
~$\Assign{\prime}{M}(x)$ (isto é, é ``fechado sob o sucessor''), então
contém tudo no !!{domain} (isto é, $X = \Domain{M}$).

O axioma de indução garante que qualquer !!{structure} que o satisfaça
contenha apenas os !!{element}s de~$\Domain{M}$ cuja existência os axiomas
exigem, isto é, os valores de~$\num{n}$ para $n \in \Nat$. Um modelo de
$\Th{PA^2}$ não contém números não padrão.

\begin{thm}
\ollabel{thm:sol-pa-standard}
Se $\Sat{M}{\Th{PA^2}}$, então $\Domain{M} =
\Setabs{\Value{\num{n}}{M}}{n \in \Nat}$.
\end{thm}

\begin{proof}
Seja $N = \Setabs{\Value{\num{n}}{M}}{n \in \Nat}$, e suponha que
$\Sat{M}{\Th{PA^2}}$. É claro que, para qualquer $n \in \Nat$,
$\Value{\num{n}}{M} \in \Domain{M}$, portanto $N \subseteq \Domain{M}$.

Agora consideremos a inclusão na outra direção. Tome uma atribuição de
variáveis $s$ com $s(X) = N$. Por hipótese,
\begin{align*}
\Sat{M}{\lforall[X][((X(\Obj 0) \land \lforall[x][(X(x)
      \lif X(x'))]) \lif \lforall[x][X(x)])]}, & \text{ portanto}\\
\Sat{M}{(X(\Obj 0) \land \lforall[x][(X(x) \lif X(x'))]) \lif
  \lforall[x][X(x)]}[s]. & 
\end{align*}
Considere o antecedente dessa condicional. $\Value{\Obj{0}}{M} \in N$ e,
portanto, $\Sat{M}{X(\Obj{0})}[s]$. A segunda conjunção,
$\lforall[x][(X(x) \lif X(x'))]$, também é satisfeita. Com efeito, suponha
que $x \in N$. Pela definição de~$N$, $x = \Value{\num{n}}{M}$ para
algum~$n$. Isso dá $\Assign{\prime}{M}(x) = \Value{\num{n+1}}{M} \in N$.
Logo, $\Assign{\prime}{M}(x) \in N$.

Temos $\Sat{M}{X(\Obj 0) \land \lforall[x][(X(x) \lif X(x'))]}[s]$.
Consequentemente, $\Sat{M}{\lforall[x][X(x)]}[s]$. Mas isso significa que,
para todo $x \in \Domain{M}$, temos $x \in s(X) = N$. Assim,
$\Domain{M} \subseteq N$.
\end{proof}

\begin{cor}
\ollabel{cor:sol-pa-categorical}
Quaisquer dois modelos de~$\Th{PA^2}$ são isomorfos.
\end{cor}

\begin{proof}
Por \olref{thm:sol-pa-standard}, o domínio de qualquer modelo de
$\Th{PA^2}$ é esgotado por $\Value{\num{n}}{M}$. Todo modelo desse tipo
também é modelo de~$\Th{Q}$. Por
\olref[mod][mar][stm]{prop:thq-standard}, todo modelo desse tipo é padrão,
isto é, isomorfo a~$\Struct{N}$.
\end{proof}

Acima definimos $\Th{PA^2}$ como a teoria que contém os oito primeiros
axiomas aritméticos mais o axioma de indução de segunda ordem. Na verdade,
graças ao poder expressivo da lógica de segunda ordem, apenas os
\emph{dois primeiros} axiomas aritméticos, juntamente com a indução, são
necessários para a aritmética de Peano de segunda ordem.

\begin{prop}\ollabel{prop:sol-pa-definable}
Seja $\Th{PA^{2\dagger}}$ a teoria de segunda ordem que contém os dois
primeiros axiomas aritméticos (os axiomas do sucessor) e o axioma de indução
de segunda ordem. Então $\le$, $+$ e $\times$ são definíveis em
$\Th{PA^{2\dagger}}$.
\end{prop}

\begin{proof}
Para mostrar que $\le$ é definível, precisamos encontrar uma
fórmula~$!A_\le(x, y)$ tal que $\Sat{N}{!A_\le(\num n, \num m)}$ sse
$n \le m$. Considere a fórmula
\[
!B(x, Y) \ident Y(x) \land \lforall[y][(Y(y) \lif Y(y'))]
\]
É claro que $!B(\num n, Y)$ é satisfeita por um conjunto
$Y \subseteq \Nat$ sse $\Setabs{m}{n \le m} \subseteq Y$; portanto,
podemos tomar $!A_\le(x, y) \ident
\lforall[Y][(!B(x, Y) \lif Y(y))]$.

Para ver que a adição é definível, observe que $k+l=m$ sse existe uma função
$u$ tal que $u(0)=k$, $u(n')=u(n)'$ para todo $n$, e $m = u(l)$. Podemos
usar essa equivalência para definir a adição em $\Th{PA^{2\dagger}}$ pela
fórmula seguinte:
\[
!A_+(x,y,z) \ident \lexists[u][(u(\Obj 0)=x \land \lforall[w][u(x')=u
 (x)'] \land u(y)=z)]
\] 
Deve estar claro que $\Sat{N}{!A_+(\num k, \num l, \num m)}$ sse
$k+l=m$.
\end{proof}

\begin{prob}
Complete a prova de \olref[sol][met][spa]{prop:sol-pa-definable}.
\end{prob}

\end{document}
